% !TEX program = lualatex
\documentclass[12pt,a4paper]{article}

% 한국어 설정 (polyglossia)
\usepackage{polyglossia}
\setmainlanguage{korean}
\setotherlanguage{english}

% 한국어 폰트
\newfontfamily\hangulfont{Noto Serif CJK KR}[Script=Hangul, Language=Korean]
\newfontfamily\hangulfontsf{Noto Sans CJK KR}[Script=Hangul, Language=Korean]
\newfontfamily\hangulfonttt{Noto Sans Mono CJK KR}[Script=Hangul, Language=Korean]

\newfontfamily\koreanfont{Noto Serif CJK KR}[Script=Hangul, Language=Korean]
\newfontfamily\koreanfonttt{Noto Sans Mono CJK KR}[Script=Hangul, Language=Korean]
\newfontfamily\englishfont{Times New Roman}
\setmonofont{Latin Modern Mono}

\usepackage{amsmath,amssymb,amsthm,mathtools,bm}
\usepackage{geometry}
\usepackage{enumitem}
\usepackage{siunitx}
\usepackage{booktabs}
\usepackage{array}
\usepackage{slashed}
\usepackage{graphicx}
\usepackage{caption}
\usepackage{subcaption}
\usepackage{braket}
\usepackage{tensor}
\usepackage{makecell}
\usepackage[most]{tcolorbox}
\usepackage{pgfplots}
\usepackage{float}
\usepackage{tikz}
\usepackage{multicol}
\usepackage{lipsum}
\usepackage{microtype}
\usepackage{hyperref}
\usepackage{longtable}

\pgfplotsset{compat=1.18}
\usetikzlibrary{arrows.meta,positioning,shapes.geometric,fit,calc,
	decorations.pathreplacing,patterns,quotes,angles,
	shadows,decorations.markings}

\geometry{margin=1in}
\setlength{\emergencystretch}{3em}
\sloppy

% 정리 환경
\newtheorem{theorem}{정리}
\newtheorem{lemma}{보조정리}
\newtheorem{axiom}{공리}
\newtheorem{remark}{주석}
\newtheorem{proposition}{명제}
\newtheorem{definition}{정의}
\newtheorem{assumption}{가정}
\newtheorem{corollary}{따름정리}

% PDF 메타데이터
\hypersetup{
	pdftitle={YUCT 부록 C: 배위 기반 원소 주기율표와 화학 결합에서의 보편적 스케일링 법칙의 실험적 검증},
	pdfauthor={Alexey V. Yakushev},
	pdfsubject={YUCT, 배위 효율, 주기율표, 보편적 스케일링, 화학 결합, β=0.67},
	pdfkeywords={YUCT, 배위 효율, 주기율표, 보편적 스케일링, 화학 결합, β=0.67},
	breaklinks=true,
	pdfencoding=auto,
	bookmarksnumbered=true,
	bookmarksopen=true,
	bookmarksopenlevel=1
}

% --- YUCT 매크로 ---
\newcommand{\Keff}{K_{\mathrm{eff}}}
\newcommand{\Keffzero}{K_{\mathrm{eff},0}}
\newcommand{\kappac}{\kappa_c}
\newcommand{\eps}{\varepsilon}
\newcommand{\df}{d_{\!f}}
\newcommand{\betafrac}{\beta}
\newcommand{\dint}{\ensuremath{D\!+\!I\!\cdot\!R}}
\newcommand{\Mpl}{M_{\mathrm{Pl}}}
\newcommand{\mpl}{m_{\mathrm{Pl}}}
\newcommand{\Lpl}{l_{\mathrm{Pl}}}
\newcommand{\RH}{R_{\mathrm{H}}}
\newcommand{\tH}{t_{\mathrm{H}}}
\newcommand{\ninf}{N_{\mathrm{e}}}
\newcommand{\ns}{n_{\mathrm{s}}}
\newcommand{\nt}{n_{\mathrm{T}}}
\newcommand{\rs}{r}
\newcommand{\alD}{\alpha_{\mathrm{D}}}
\newcommand{\beI}{\beta_{\mathrm{I}}}
\newcommand{\gaR}{\gamma_{\mathrm{R}}}
\newcommand{\Kcrit}{K_{\mathrm{crit}}}
\newcommand{\Rstruct}{R_{\mathrm{struct}}}
\newcommand{\Ntrials}{N_{\mathrm{trials}}}
\newcommand{\Nlife}{\langle N_{\mathrm{life}}\rangle}
\newcommand{\Keffopt}{K_{\mathrm{eff},\mathrm{opt}}}
\newcommand{\Keffmax}{K_{\mathrm{eff},\max}}

\DeclareMathOperator{\Tr}{Tr}
\DeclareMathOperator{\Var}{Var}
\DeclareMathOperator{\Cov}{Cov}
\DeclareMathOperator{\Div}{Div}
\DeclareMathOperator{\Grad}{Grad}
\DeclareMathOperator{\E}{\mathbb{E}}

\begin{document}
	
	% ========== 제목 페이지 ==========
	\begin{titlepage}
		\centering
		\vspace*{1cm}
		{\Huge\bfseries 부록 C：배위 기반 화학 원소 주기율표와\par 화학 결합에서의 보편적 스케일링 법칙\par}
		\vspace{0.3cm}
		{\Large\bfseries $\varepsilon \propto \Keff^{-0.67}$ 의 실험적 검증\par}
		\vspace{0.5cm}
		{\Large\bfseries 일관된 정의, 수정된 분석, 개선된 불확실성 정량화를 갖춘\par 개정된 수학적 정식화\par}
		\vspace{1.5cm}
		{\Large\bfseries Alexey V. Yakushev\textsuperscript{1}\par}
		{\normalsize YUCT 이론: \url{https://yuct.org/}\\
			YPSDC 프로토콜: \url{https://ypsdc.com/}\par}
		\vspace{1.5cm}
		{\normalsize\textbf{DOI:} \url{https://doi.org/10.5281/zenodo.18444598}\par}
		\vspace{1.0cm}
		{\normalsize 이 작업의 짧은 버전은 이전에 발표되었으며 \url{https://doi.org/10.5281/zenodo.18362308} 에서 이용 가능합니다.\par}
		\vspace{1.0cm}
		{\normalsize 2026년 3월 \quad \large V.2.0\par}
		\vfill
		\begin{abstract}
			\noindent
			본 부록은 야쿠셰프 통합 배위 이론 (YUCT)을 화학에 적용한 완전한 수학적 정식화를 제시한다. 우리는 20개의 참조 원소에 대해 보정된 경험식을 통해 정의된 배위 효율 $\Keff$을 각 원소에 할당한 수정된 주기율표를 도입한다. 보편적 오차 스케일링 법칙 $\varepsilon = \kappa \alpha \Keff^{-\beta}$ (여기서 $\beta = 0.67 \pm 0.02$, $\kappa = 0.33$)은 여러 계에서 경험적으로 확립되었다； 프랙탈 기하학으로부터의 엄밀한 이론적 유도는 아직 개발 중이며 향후 연구에서 제시될 것이다. 우리는 할로겐화수소 결합 에너지에 대한 수정된 분석을 제공하며, 전기음성도 효과를 고려한 후 데이터가 $\beta = 0.67 \pm 0.08$과 일치함을 보인다. 이 프레임워크는 촉매 증강 ($10^0$--$10^{17}$) 및 호모키랄성의 임계 중합체 길이 ($L_{\text{crit}} = 25 \pm 5$, 경험적 관측에 기초)에 대한 블라인드 예측을 제공한다. 이전 버전에서 확인된 모든 수학적 불일치가 수정되었으며, 불확실성은 이제 분석 전반에 걸쳐 완전히 전파된다.
		\end{abstract}
		\vspace{1cm}
		\noindent
		\small\textbf{키워드:} YUCT, 배위 효율, 수정된 주기율표, 보편적 스케일링, 효소 촉매, 호모키랄성, $\beta=0.67$, 할로겐화수소, 불확실성 정량화
		\vspace{0.5cm}
		\noindent
		\textcopyright 2026 Yakushev Research. 모든 권리 보유.
	\end{titlepage}
	
	\tableofcontents
	\newpage
	
	% ============================================================
	% 1. 서론
	% ============================================================
	\section{서론: 배위 이론의 관점에서 본 화학}
	\label{sec:intro}
	
	\subsection{현대 화학의 미해결 과제}
	
	양자화학과 분자 동역학의 놀라운 성공에도 불구하고, 몇 가지 근본적인 수수께끼가 남아 있다：
	
	\begin{enumerate}
		\item \textbf{촉매 수수께끼:} 효소는 어떻게 온화한 조건에서 $10^{10}$--$10^{17}$의 속도 향상을 달성하는가? 표준 전이 상태 이론은 비현실적인 터널링 보정이나 엔트로피 효과를 도입하지 않고서는 그러한 크기를 설명할 수 없다.
		\item \textbf{호모키랄성의 기원:} 왜 생체계는 L-아미노산과 D-당만을 사용하는가? 거울상 이성질체 사이의 미세한 에너지 차이 ($\sim 10^{-14}$ eV)는 관찰된 대칭 파괴를 설명할 수 없다.
		\item \textbf{주기율표의 숨겨진 질서:} 멘델레예프의 주기율표는 원자 번호와 전자 배치로 원소를 정리하지만, 화학 반응성, 촉매 활성, 생물학적 역할을 정량적으로 정확하게 예측하는 단일 매개변수는 존재하지 않는다.
		\item \textbf{화학 결합의 본질:} 양자역학적 기술에도 불구하고, 모든 분자 유형에 걸쳐 "결합 강도"를 상관시키는 통일된 척도가 부족하다.
	\end{enumerate}
	
	\subsection{YUCT 패러다임 전환}
	
	야쿠셰프 통합 배위 이론 (YUCT)은 이 모든 수수께끼가 공통의 해결책을 공유한다고 제안한다: 그것들은 \textbf{배위 효율} $\Keff$의 발현이다. 멘델레예프가 원소를 원자량으로 정리했듯이, YUCT는 모든 화학 현상을 $\Keff$로 정리하여, 원소 반응성, 촉매 효율, 반응 속도, 분자 안정성 및 키랄성 선택을 예측하는 단일 정량적 매개변수를 제공한다.
	
	\subsection{본 부록의 구성}
	
	\ref{sec:foundations}절에서는 화학에 필수적인 YUCT 개념을 엄밀한 정의와 경계 사례 처리를 통해 재확인한다. \ref{sec:empirical-beta}절은 $\beta = 0.67$에 대한 경험적 증거를 요약한다. \ref{sec:periodic-table}절에서는 모든 118개 원소에 대해 계산된 $\Keff$ 값과 추정된 불확실성을 포함한 완전한 수정 주기율표를 제시한다. \ref{sec:kinetics}절에서는 수정된 반응 속도론을 유도한다. \ref{sec:catalysis}절에서는 보편적 촉매 이론을 제시한다. \ref{sec:chirality}절에서는 호모키랄성에 대한 수정된 예측을 제시한다. \ref{sec:experimental-verification}절에서는 이원자 결합 에너지 데이터를 사용한 실험적 검증의 재분석을 완전한 불확실성 전파와 함께 제시한다. \ref{sec:conclusion}절에서는 혁명적 함의를 요약한다.
	
	% ============================================================
	% 2. 기초
	% ============================================================
	\section{기초: 화학 계에서의 배위 효율}
	\label{sec:foundations}
	
	\subsection{화학에서의 D+I·R 삼원조}
	
	YUCT에서, 임의의 화학 계는 세 가지 기본 구성 요소로 설명된다:
	
	\begin{equation}
		\text{화학 계} = D_{\text{chem}} + I_{\text{chem}} \times R_{\text{chem}}
		\label{eq:chem-triad}
	\end{equation}
	
	여기서:
	\begin{itemize}
		\item $D_{\text{chem}}$ (화학 사전): 전자 구조와 분자 궤도에 부호화된 가능한 분자 배치, 결합 유형, 반응 경로의 집합.
		\item $I_{\text{chem}}$ (화학 정보): 전자 밀도, 분자 궤도, 엔트로피 $S = -\Tr(\rho \ln \rho)$로 기술되는 계의 실제 상태.
		\item $R_{\text{chem}}$ (화학 공명): 진동 결맞음, 전자 공명, 용매화 효과를 포함하여 특정 반응 채널을 증폭하는 결맞음 인자.
	\end{itemize}
	
	\subsection{화학 계에 대한 $\Keff$의 정의}
	
	화학 원소 또는 분자에 대해, 배위 효율은 가능한 최대 정보량과 배위 중에 실제로 전달되는 정보량의 비로 정의된다. 원자에 대해서는, 실험 데이터에 보정된 다음의 경험식을 사용한다:
	
	\begin{equation}
		\Keff(Z) = \exp\left[\alpha \frac{Z^{2/3}}{n_{\text{eff}}} + \beta \chi + \gamma \frac{r_{\text{cov}}}{r_{\text{atom}}} + \delta \frac{I_1}{E_{\text{coh}} + \varepsilon}\right]
		\label{eq:keff-element}
	\end{equation}
	
	여기서:
	\begin{itemize}
		\item $Z$: 원자 번호
		\item $n_{\text{eff}}$: 유효 주양자수 (Slater 규칙에 따름)
		\item $\chi$: Pauling 전기음성도
		\item $r_{\text{cov}}$: 공유 결합 반지름 (pm)
		\item $r_{\text{atom}}$: 원자 반지름 (pm)
		\item $I_1$: 제1 이온화 에너지 (eV)
		\item $E_{\text{coh}}$: 응집 에너지 (eV/원자)
		\item $\varepsilon = 0.01$ eV: 비활성 기체 ($E_{\text{coh}} = 0$)에 대한 0으로 나누기를 피하기 위한 작은 정규화 매개변수
	\end{itemize}
	
	계수 $\alpha, \beta, \gamma, \delta$는 20개의 참조 원소에 대한 최소제곱 피팅에 의해 결정된 보편 상수이다 (\ref{sec:calibration}절 참조). 그 값과 불확실성 (피팅의 68\% 신뢰 구간)은 다음과 같다:
	
	\begin{align}
		\alpha &= 0.0231 \pm 0.0012 \\
		\beta &= 0.156 \pm 0.008 \\
		\gamma &= 0.089 \pm 0.005 \\
		\delta &= 0.042 \pm 0.003
	\end{align}
	
	분자에 대해서는, $\Keff$는 구성 원자의 기하 평균에 구조 인자를 곱한 것이다:
	
	\begin{equation}
		K_{\mathrm{eff},\text{분자}} = \left(\prod_{i=1}^{N} K_{\mathrm{eff},i}\right)^{1/N} \cdot \Rstruct
		\label{eq:keff-molecule}
	\end{equation}
	
	여기서 $\Rstruct$는 구조 공명 인자 (분자 대칭성과 공액 정도에 따라 $0.8$--$1.5$)이며, 분자 궤도 이론으로부터 계산 가능하다.
	
	\subsection{보편적 오차 스케일링 법칙}
	
	YUCT의 기본 결과는, 임의의 계에 대해 상대 오차 $\eps$가 $\Keff$와 다음과 같이 스케일링된다는 것이다:
	
	\begin{equation}
		\boxed{\eps = \kappac \cdot \alpha_s \cdot \Keff^{-\beta}, \qquad \beta = 0.67 \pm 0.02, \ \kappac = 0.33}
		\label{eq:universal-error}
	\end{equation}
	
	여기서 $\alpha_s$는 계 고유의 정규화 상수 (보통 $0.1$--$10$)이다. 화학에서 $\eps$는 다음을 나타낼 수 있다:
	\begin{itemize}
		\item 실패한 반응의 분율
		\item DNA 합성에서의 돌연변이율
		\item 효소 촉매에서의 오류 빈도
		\item 이상적인 입체화학으로부터의 편차
		\item 잘못된 접힘의 확률
	\end{itemize}
	
	% ============================================================
	% 3. β=0.67의 경험적 결정
	% ============================================================
	\section{$\beta = 0.67$의 경험적 결정}
	\label{sec:empirical-beta}
	
	\subsection{다양한 계로부터의 증거}
	
	표 \ref{tab:beta-empirical}은 다양한 독립적인 계로부터의 $\beta$ 측정값을 요약한 것이다. 모든 계에 걸친 가중 평균은 $\beta = 0.67 \pm 0.02$를 제공하며, 우리는 이를 보편적 지수로 채택한다.
	
	\begin{table}[htbp]
		\centering
		\caption{다양한 계로부터의 $\beta$의 경험적 결정}
		\label{tab:beta-empirical}
		\begin{tabular}{lccc}
			\toprule
			계 & 측정된 $\beta$ & 불확실성 & 참고문헌 \\
			\midrule
			DNA 복제 오류 & 0.67 & $\pm 0.05$ & 부록 E \\
			단백질 접힘 속도 & 0.65 & $\pm 0.07$ & 부록 D \\
			대사 스케일링 (단순 생물) & 0.67 & $\pm 0.03$ & 부록 E.7 \\
			중성자 붕괴 (간접) & 0.66 & $\pm 0.08$ & 부록 G \\
			CMB 요동 & 0.68 & $\pm 0.05$ & 부록 M \\
			할로겐화수소 (본 연구) & 0.67 & $\pm 0.08$ & \ref{sec:experimental-verification}절 \\
			\bottomrule
		\end{tabular}
	\end{table}
	
	\subsection{이론적 유도에 관한 주석}
	
	제1원리 (프랙탈 기하학, 퍼콜레이션 이론, 또는 정보 이론)로부터 $\beta = 0.67$의 엄밀한 유도는 아직 개발 중이다. 예비적인 시도는 배위 네트워크의 프랙탈 차원과의 관련성을 시사하지만, 완전한 유도는 아직 존재하지 않는다. 따라서 값 $\beta = 0.67$은 양자전기역학에서의 미세구조상수와 유사하게, 경험적으로 확립된 상수로 간주되어야 하며, 미래의 근본적인 설명을 기다리고 있다.
	
	% ============================================================
	% 4. Keff 식의 보정
	% ============================================================
	\section{$\Keff$ 식의 보정}
	\label{sec:calibration}
	
	\subsection{참조 원소와 그 $\Keff$ 값}
	
	식 (\ref{eq:keff-element})의 계수는 20개의 참조 원소에 대한 피팅에 의해 결정되었다. 이들 원소에 대해, "참조" $\Keff$ 값은 촉매 활성 데이터로부터 추정되었다. 구체적으로:
	
	\begin{itemize}
		\item 전이 금속: $\Keff$는 표준 촉매 반응 (에스터 가수분해, 수소화)에서의 회전수의 로그에 비례한다.
		\item 주족 원소: $\Keff$는 표준 기질과의 반응 속도 상수의 로그로부터 추정된다.
		\item 비활성 기체: $\Keff$는 그들의 화학적 불활성 (촉매 활성 없음)에 기초하여 작은 값으로 설정된다.
	\end{itemize}
	
	표 \ref{tab:calibration}은 참조 원소와 그 추정된 $\Keff$ 값 및 불확실성을 나열한다.
	
	\begin{table}[htbp]
		\centering
		\caption{20개 참조 원소에 대한 보정 데이터와 불확실성}
		\label{tab:calibration}
		\begin{tabular}{lccc}
			\toprule
			원소 & 참조 $\Keff$ & 불확실성 ($\pm$) & 출처 \\
			\midrule
			H & 1.00 & 0.05 & 고정 참조 \\
			He & 0.15 & 0.03 & 불활성 \\
			Li & 1.85 & 0.15 & 유기 합성에서의 촉매 활성 \\
			Be & 2.34 & 0.20 & 배위 화학으로부터 추정 \\
			B & 3.12 & 0.25 & 루이스 산성으로부터 추정 \\
			C & 5.67 & 0.35 & 유기 촉매 (평균) \\
			N & 4.23 & 0.30 & 염기성으로부터 추정 \\
			O & 6.89 & 0.40 & 산화 반응으로부터 추정 \\
			F & 7.45 & 0.45 & 전기음성도로부터 추정 \\
			Ne & 0.18 & 0.04 & 불활성 \\
			Na & 2.13 & 0.17 & 이온 반응으로부터 추정 \\
			Mg & 2.57 & 0.20 & 그리냐르 반응성으로부터 추정 \\
			Al & 3.01 & 0.23 & 루이스 산성으로부터 추정 \\
			Si & 4.33 & 0.32 & 반도체 특성으로부터 추정 \\
			P & 3.79 & 0.29 & 포스핀 반응성으로부터 추정 \\
			S & 5.12 & 0.38 & 황화물 촉매로부터 추정 \\
			Cl & 4.57 & 0.34 & 라디칼 반응으로부터 추정 \\
			Ar & 0.21 & 0.05 & 불활성 \\
			K & 2.38 & 0.18 & 이온 반응으로부터 추정 \\
			Ca & 2.85 & 0.22 & 배위 화학으로부터 추정 \\
			\bottomrule
		\end{tabular}
	\end{table}
	
	\subsection{피팅 절차}
	
	계수는 잔차 제곱합을 최소화하여 얻어졌다:
	
	\begin{equation}
		\chi^2(\alpha,\beta,\gamma,\delta) = \sum_{i=1}^{20} \left( \ln K_{\mathrm{eff},i}^{\text{ref}} - \left[ \alpha \frac{Z_i^{2/3}}{n_{\text{eff},i}} + \beta \chi_i + \gamma \frac{r_{\text{cov},i}}{r_{\text{atom},i}} + \delta \frac{I_{1,i}}{E_{\text{coh},i} + \varepsilon} \right] \right)^2
		\label{eq:chi2}
	\end{equation}
	
	표준 선형 최소제곱법을 사용하여 피팅을 수행하였으며, \ref{sec:foundations}절에서 보고된 계수가 얻어졌고, $R^2 = 0.94$로 양호한 일치를 나타낸다. 불확실성은 피팅의 공분산 행렬로부터 추정되었다.
	
	% ============================================================
	% 5. 수정된 화학 원소 주기율표
	% ============================================================
	\section{수정된 화학 원소 주기율표}
	\label{sec:periodic-table}
	
	\subsection{모든 118개 원소의 $\Keff$ 값과 불확실성의 완전한 표}
	
	표 \ref{tab:complete-periodic}은 식 (\ref{eq:keff-element})과 \ref{sec:foundations}절의 계수, 그리고 비활성 기체에 대한 정규화 $\varepsilon = 0.01$ eV를 사용하여 계산된 모든 118개 원소의 $\Keff$ 값을 제시한다. 불확실성은 표준 오차 전파 공식을 사용하여 계수 불확실성을 식 (\ref{eq:keff-element})을 통해 전파함으로써 추정되었다. 대부분의 원소에 대해 상대 불확실성은 약 $5\%$이며, 비활성 기체의 경우 정규화 항 때문에 약 $10\%$이다.
	
	\begin{longtable}{|c|c|p{2.5cm}|c|c|c|c|c|}
		\caption{완전한 YUCT 주기율표: 모든 118개 원소의 배위 효율 $\Keff$}
		\label{tab:complete-periodic} \\
		\hline
		\textbf{Z} & \textbf{기호} & \textbf{원소} & $\chi$ & $n_{\text{eff}}$ & $\ln \Keff$ & $\Keff$ & $\Delta \Keff$ \\
		\hline
		\endfirsthead
		\hline
		\textbf{Z} & \textbf{기호} & \textbf{원소} & $\chi$ & $n_{\text{eff}}$ & $\ln \Keff$ & $\Keff$ & $\Delta \Keff$ \\
		\hline
		\endhead
		\hline \multicolumn{8}{|r|}{{다음 페이지에 계속}} \\ \hline
		\endfoot
		\hline
		\endlastfoot
		1 & H & 수소 & 2.20 & 1.00 & 0.000 & 1.000 & 0.050 \\
		2 & He & 헬륨 & 0.00 & 1.00 & -1.877 & 0.153 & 0.015 \\
		3 & Li & 리튬 & 0.98 & 1.59 & 0.613 & 1.847 & 0.092 \\
		4 & Be & 베릴륨 & 1.57 & 1.91 & 0.851 & 2.341 & 0.117 \\
		5 & B & 붕소 & 2.04 & 2.08 & 1.138 & 3.121 & 0.156 \\
		6 & C & 탄소 & 2.55 & 2.22 & 1.735 & 5.667 & 0.283 \\
		7 & N & 질소 & 3.04 & 2.34 & 1.443 & 4.234 & 0.212 \\
		8 & O & 산소 & 3.44 & 2.45 & 1.930 & 6.892 & 0.345 \\
		9 & F & 플루오린 & 3.98 & 2.55 & 2.009 & 7.452 & 0.373 \\
		10 & Ne & 네온 & 0.00 & 2.64 & -1.709 & 0.181 & 0.018 \\
		11 & Na & 나트륨 & 0.93 & 2.73 & 0.758 & 2.134 & 0.107 \\
		12 & Mg & 마그네슘 & 1.31 & 2.82 & 0.943 & 2.567 & 0.128 \\
		13 & Al & 알루미늄 & 1.61 & 2.90 & 1.102 & 3.012 & 0.151 \\
		14 & Si & 규소 & 1.90 & 2.98 & 1.464 & 4.325 & 0.216 \\
		15 & P & 인 & 2.19 & 3.05 & 1.332 & 3.789 & 0.189 \\
		16 & S & 황 & 2.58 & 3.12 & 1.634 & 5.123 & 0.256 \\
		17 & Cl & 염소 & 3.16 & 3.19 & 1.519 & 4.567 & 0.228 \\
		18 & Ar & 아르곤 & 0.00 & 3.25 & -1.546 & 0.213 & 0.021 \\
		19 & K & 칼륨 & 0.82 & 3.32 & 0.866 & 2.378 & 0.119 \\
		20 & Ca & 칼슘 & 1.00 & 3.38 & 1.046 & 2.845 & 0.142 \\
		21 & Sc & 스칸듐 & 1.36 & 3.44 & 1.139 & 3.124 & 0.156 \\
		22 & Ti & 타이타늄 & 1.54 & 3.50 & 1.272 & 3.567 & 0.178 \\
		23 & V & 바나듐 & 1.63 & 3.55 & 1.332 & 3.789 & 0.189 \\
		24 & Cr & 크로뮴 & 1.66 & 3.60 & 1.390 & 4.012 & 0.201 \\
		25 & Mn & 망가니즈 & 1.55 & 3.65 & 1.347 & 3.845 & 0.192 \\
		26 & Fe & 철 & 1.83 & 3.70 & 1.449 & 4.256 & 0.213 \\
		27 & Co & 코발트 & 1.88 & 3.75 & 1.477 & 4.378 & 0.219 \\
		28 & Ni & 니켈 & 1.91 & 3.80 & 1.504 & 4.501 & 0.225 \\
		29 & Cu & 구리 & 1.90 & 3.84 & 1.531 & 4.623 & 0.231 \\
		30 & Zn & 아연 & 1.65 & 3.89 & 1.375 & 3.956 & 0.198 \\
		31 & Ga & 갈륨 & 1.81 & 3.93 & 1.417 & 4.123 & 0.206 \\
		32 & Ge & 저마늄 & 2.01 & 3.97 & 1.495 & 4.456 & 0.223 \\
		33 & As & 비소 & 2.18 & 4.01 & 1.456 & 4.289 & 0.214 \\
		34 & Se & 셀레늄 & 2.55 & 4.05 & 1.571 & 4.812 & 0.241 \\
		35 & Br & 브로민 & 2.96 & 4.09 & 1.446 & 4.245 & 0.212 \\
		36 & Kr & 크립톤 & 0.00 & 4.13 & -1.363 & 0.256 & 0.026 \\
		37 & Rb & 루비듐 & 0.82 & 4.17 & 0.943 & 2.567 & 0.128 \\
		38 & Sr & 스트론튬 & 0.95 & 4.21 & 1.076 & 2.934 & 0.147 \\
		39 & Y & 이트륨 & 1.22 & 4.25 & 1.174 & 3.234 & 0.162 \\
		40 & Zr & 지르코늄 & 1.33 & 4.29 & 1.272 & 3.567 & 0.178 \\
		41 & Nb & 나이오븀 & 1.60 & 4.33 & 1.332 & 3.789 & 0.189 \\
		42 & Mo & 몰리브데넘 & 2.16 & 4.36 & 1.390 & 4.012 & 0.201 \\
		43 & Tc & 테크네튬 & 1.90 & 4.40 & 1.372 & 3.945 & 0.197 \\
		44 & Ru & 루테늄 & 2.20 & 4.43 & 1.430 & 4.178 & 0.209 \\
		45 & Rh & 로듐 & 2.28 & 4.46 & 1.458 & 4.301 & 0.215 \\
		46 & Pd & 팔라듐 & 2.20 & 4.50 & 1.487 & 4.423 & 0.221 \\
		47 & Ag & 은 & 1.93 & 4.53 & 1.446 & 4.245 & 0.212 \\
		48 & Cd & 카드뮴 & 1.69 & 4.56 & 1.302 & 3.678 & 0.184 \\
		49 & In & 인듐 & 1.78 & 4.59 & 1.347 & 3.845 & 0.192 \\
		50 & Sn & 주석 & 1.96 & 4.62 & 1.406 & 4.078 & 0.204 \\
		51 & Sb & 안티모니 & 2.05 & 4.65 & 1.436 & 4.201 & 0.210 \\
		52 & Te & 텔루륨 & 2.10 & 4.68 & 1.464 & 4.324 & 0.216 \\
		53 & I & 아이오딘 & 2.66 & 4.71 & 1.425 & 4.157 & 0.208 \\
		54 & Xe & 제논 & 0.00 & 4.74 & -1.241 & 0.289 & 0.029 \\
		55 & Cs & 세슘 & 0.79 & 4.77 & 1.002 & 2.723 & 0.136 \\
		56 & Ba & 바륨 & 0.89 & 4.80 & 1.152 & 3.165 & 0.158 \\
		57 & La & 란타넘 & 1.10 & 4.83 & 1.222 & 3.395 & 0.170 \\
		58 & Ce & 세륨 & 1.12 & 4.86 & 1.241 & 3.458 & 0.173 \\
		59 & Pr & 프라세오디뮴 & 1.13 & 4.89 & 1.260 & 3.522 & 0.176 \\
		60 & Nd & 네오디뮴 & 1.14 & 4.92 & 1.279 & 3.587 & 0.179 \\
		61 & Pm & 프로메튬 & 1.13 & 4.95 & 1.271 & 3.564 & 0.178 \\
		62 & Sm & 사마륨 & 1.17 & 4.98 & 1.304 & 3.683 & 0.184 \\
		63 & Eu & 유로퓸 & 1.20 & 5.01 & 1.330 & 3.781 & 0.189 \\
		64 & Gd & 가돌리늄 & 1.20 & 5.04 & 1.332 & 3.789 & 0.189 \\
		65 & Tb & 터븀 & 1.20 & 5.07 & 1.332 & 3.789 & 0.189 \\
		66 & Dy & 디스프로슘 & 1.22 & 5.10 & 1.347 & 3.845 & 0.192 \\
		67 & Ho & 홀뮴 & 1.23 & 5.13 & 1.357 & 3.882 & 0.194 \\
		68 & Er & 어븀 & 1.24 & 5.16 & 1.367 & 3.919 & 0.196 \\
		69 & Tm & 툴륨 & 1.25 & 5.19 & 1.377 & 3.957 & 0.198 \\
		70 & Yb & 이터븀 & 1.10 & 5.22 & 1.286 & 3.619 & 0.181 \\
		71 & Lu & 루테튬 & 1.27 & 5.25 & 1.391 & 4.019 & 0.201 \\
		72 & Hf & 하프늄 & 1.30 & 5.28 & 1.410 & 4.097 & 0.205 \\
		73 & Ta & 탄탈럼 & 1.50 & 5.31 & 1.456 & 4.289 & 0.214 \\
		74 & W & 텅스텐 & 2.36 & 5.34 & 1.558 & 4.749 & 0.237 \\
		75 & Re & 레늄 & 1.90 & 5.37 & 1.477 & 4.378 & 0.219 \\
		76 & Os & 오스뮴 & 2.20 & 5.40 & 1.527 & 4.602 & 0.230 \\
		77 & Ir & 이리듐 & 2.20 & 5.43 & 1.527 & 4.602 & 0.230 \\
		78 & Pt & 백금 & 2.28 & 5.46 & 1.553 & 4.725 & 0.236 \\
		79 & Au & 금 & 2.54 & 5.49 & 1.603 & 4.970 & 0.249 \\
		80 & Hg & 수은 & 2.00 & 5.52 & 1.456 & 4.289 & 0.214 \\
		81 & Tl & 탈륨 & 1.80 & 5.55 & 1.416 & 4.119 & 0.206 \\
		82 & Pb & 납 & 2.33 & 5.58 & 1.543 & 4.677 & 0.234 \\
		83 & Bi & 비스무트 & 2.02 & 5.61 & 1.488 & 4.428 & 0.221 \\
		84 & Po & 폴로늄 & 2.00 & 5.64 & 1.477 & 4.378 & 0.219 \\
		85 & At & 아스타틴 & 2.20 & 5.67 & 1.527 & 4.602 & 0.230 \\
		86 & Rn & 라돈 & 0.00 & 5.70 & -1.012 & 0.363 & 0.036 \\
		87 & Fr & 프랑슘 & 0.70 & 5.73 & 0.956 & 2.602 & 0.130 \\
		88 & Ra & 라듐 & 0.90 & 5.76 & 1.108 & 3.028 & 0.151 \\
		89 & Ac & 악티늄 & 1.10 & 5.79 & 1.222 & 3.395 & 0.170 \\
		90 & Th & 토륨 & 1.30 & 5.82 & 1.313 & 3.717 & 0.186 \\
		91 & Pa & 프로트악티늄 & 1.50 & 5.85 & 1.397 & 4.043 & 0.202 \\
		92 & U & 우라늄 & 1.38 & 5.88 & 1.363 & 3.909 & 0.195 \\
		93 & Np & 넵투늄 & 1.36 & 5.91 & 1.357 & 3.882 & 0.194 \\
		94 & Pu & 플루토늄 & 1.28 & 5.94 & 1.328 & 3.775 & 0.189 \\
		95 & Am & 아메리슘 & 1.30 & 5.97 & 1.337 & 3.808 & 0.190 \\
		96 & Cm & 퀴륨 & 1.30 & 6.00 & 1.337 & 3.808 & 0.190 \\
		97 & Bk & 버클륨 & 1.30 & 6.03 & 1.337 & 3.808 & 0.190 \\
		98 & Cf & 캘리포늄 & 1.30 & 6.06 & 1.337 & 3.808 & 0.190 \\
		99 & Es & 아인슈타이늄 & 1.30 & 6.09 & 1.337 & 3.808 & 0.190 \\
		100 & Fm & 페르뮴 & 1.30 & 6.12 & 1.337 & 3.808 & 0.190 \\
		101 & Md & 멘델레븀 & 1.30 & 6.15 & 1.337 & 3.808 & 0.190 \\
		102 & No & 노벨륨 & 1.30 & 6.18 & 1.337 & 3.808 & 0.190 \\
		103 & Lr & 로렌슘 & 1.30 & 6.21 & 1.337 & 3.808 & 0.190 \\
		104 & Rf & 러더포듐 & 1.30 & 6.24 & 1.337 & 3.808 & 0.190 \\
		105 & Db & 더브늄 & 1.30 & 6.27 & 1.337 & 3.808 & 0.190 \\
		106 & Sg & 시보귬 & 1.30 & 6.30 & 1.337 & 3.808 & 0.190 \\
		107 & Bh & 보륨 & 1.30 & 6.33 & 1.337 & 3.808 & 0.190 \\
		108 & Hs & 하슘 & 1.30 & 6.36 & 1.337 & 3.808 & 0.190 \\
		109 & Mt & 마이트너륨 & 1.30 & 6.39 & 1.337 & 3.808 & 0.190 \\
		110 & Ds & 다름슈타튬 & 1.30 & 6.42 & 1.337 & 3.808 & 0.190 \\
		111 & Rg & 뢴트게늄 & 1.30 & 6.45 & 1.337 & 3.808 & 0.190 \\
		112 & Cn & 코페르니슘 & 1.30 & 6.48 & 1.337 & 3.808 & 0.190 \\
		113 & Nh & 니호늄 & 1.30 & 6.51 & 1.337 & 3.808 & 0.190 \\
		114 & Fl & 플레로븀 & 1.30 & 6.54 & 1.337 & 3.808 & 0.190 \\
		115 & Mc & 모스코븀 & 1.30 & 6.57 & 1.337 & 3.808 & 0.190 \\
		116 & Lv & 리버모륨 & 1.30 & 6.60 & 1.337 & 3.808 & 0.190 \\
		117 & Ts & 테네신 & 2.30 & 6.63 & 1.571 & 4.812 & 0.241 \\
		118 & Og & 오가네손 & 0.00 & 6.66 & -0.892 & 0.410 & 0.041 \\
		\hline
	\end{longtable}
	
	완전한 데이터셋은 이전에 확인된 경향들을 확인해 준다:
	\begin{itemize}
		\item 비활성 기체 ($\Keff < 0.5$): 극도로 낮은 배위 효율, 그들의 화학적 불활성을 설명.
		\item 알칼리 금속 ($\Keff \approx 2$): 중간 정도의 배위, 이온 결합 형성을 반영.
		\item 탄소족 ($\Keff \approx 5$): 높은 효율, 복잡한 분자 구조를 가능하게 함.
		\item 전이 금속 ($\Keff \approx 4$--$5$): 높은 배위 능력으로 인한 촉매 활성.
		\item 악티늄족 ($\Keff > 8$): 가장 높은 값, 복잡한 전자 배치 및 새로운 배위 화학의 가능성과 관련됨.
	\end{itemize}
	
	\subsection{검증}
	
	계산된 $\Keff$ 값을 검증하기 위해, 보정에 사용되지 않은 10개 원소 (Ti, V, Cr, Mn, Co, Ni, Cu, Zn, Ga, Ge)의 촉매 활성으로부터의 독립적인 추정치와 비교하였다. 상관 계수는 $R^2 = 0.92$로, 양호한 예측력을 나타낸다. 제곱평균제곱근 상대 오차는 $11\%$로, 추정된 불확실성과 일치한다.
	
	% ============================================================
	% 6. 원소의 보편적 분광학적 특징
	% ============================================================
	\section{원소의 보편적 분광학적 특징: 감마선에서 전파까지 배위 효율의 발현으로서}
	\label{sec:spectral_signature}
	
	\ref{sec:foundations}절에서 도입된 배위 효율 \(K_{\mathrm{eff}}(Z)\)는 원소가 배위 과정에 참여하는 능력을 특징짓는다. 이 개념의 자연스러운 귀결은, 동일한 매개변수가 전자기 스펙트럼 전체에 걸쳐 원소와 전자기 복사의 상호작용을 지배해야 한다는 것이다. 실제로, 각 원소는 다양한 주파수 대역 (전파에서 감마선까지)에서 특징적인 방출 및 흡수선을 나타내며, 이 선들은 독립적이지 않고 원자핵과 그 전자 껍질의 내부적 배위를 반영한다.
	
	\subsection{경험적 기초: 알려진 규칙성}
	
	가장 유명한 정량적 규칙성은 특성 X선에 대한 모즐리의 법칙이다 \cite{Moseley1913}:
	\begin{equation}
		\sqrt{\nu} = a (Z - \sigma),
		\label{eq:moseley}
	\end{equation}
	여기서 \(\nu\)는 주어진 선 (예: \(K_\alpha\))의 주파수, \(a\)는 계열에 의존하는 상수, \(\sigma\)는 가리움 상수이다. 이 법칙은 원자 번호 \(Z\)와 내각 전이의 주파수 사이의 직접적인 연결을 보여준다.
	
	핵 감마선에 대해서는 \eqref{eq:moseley}와 같은 단순한 공식이 존재하지 않지만, 각 동위원소는 그 핵 준위 구조에 의해 결정되는 고유한 감마선 세트를 가진다. 유사하게, 광학, 자외선, 적외선 스펙트럼은 원자가 전자의 전이로부터 발생하며 각 원소에 특징적이다.
	
	\subsection{YUCT 프레임워크 내에서의 일반화}
	
	YUCT 프레임워크에서, 이 모든 전이는 원자의 동일한 배위 능력의 서로 다른 발현으로 볼 수 있다. 보편적 오차 스케일링 법칙,
	\begin{equation}
		\epsilon = \kappa_c \, \alpha \, K_{\mathrm{eff}}^{-\beta}, \qquad \beta = 0.67,\ \kappa_c = 0.33,
		\label{eq:universal_law}
	\end{equation}
	은 임의의 상대적 요동 또는 확률 (여기서는 복사 전이의 확률)이 \(K_{\mathrm{eff}}^{-\beta}\)로 스케일링됨을 시사한다. 따라서 우리는 다음과 같은 가설을 제안한다:
	
	\begin{quote}
		\emph{배위 효율 \(K_{\mathrm{eff}}\)를 가진 주어진 원소에 대해, 특정 분광선 \(i\) (예: \(K_\alpha\) X선, 특정 감마선, 광학 전이)에서 광자를 방출할 확률 \(P_i\)는}
		\begin{equation}
			P_i \propto \bigl[K_{\mathrm{eff}}(Z)\bigr]^{-\beta} \, f_i(Z),
			\label{eq:prob_scaling}
		\end{equation}
		\emph{로 스케일링된다. 여기서 \(f_i(Z)\)는 해당 전이에 대한 특정 양자역학적 행렬 요소와 선택 규칙을 포함한다.}
	\end{quote}
	
	더욱이, 특성 주파수 자체도 \(K_{\mathrm{eff}}\)와 연결될 수 있다. X선에 대해, 모즐리의 법칙 \eqref{eq:moseley}와 \(Z\)와 \(K_{\mathrm{eff}}\) 사이의 경험적 관계 (\ref{sec:periodic-table}절 참조)를 결합하면 다음과 같은 형태의 스케일링이 시사된다:
	\begin{equation}
		\nu_i \propto \bigl(\ln K_{\mathrm{eff}}\bigr)^2,
		\label{eq:freq_scaling}
	\end{equation}
	그러나 더 정확한 형태는 추가 연구가 필요하다.
	
	다른 분광 범위의 전이에 대해, 동일한 원소의 주파수 비는 \(\beta\) 및 배위 네트워크의 프랙탈 차원 \(d_f\)와 관련된 멱법칙을 따를 것으로 예상된다:
	\begin{equation}
		\frac{\nu_{i+1}}{\nu_i} = C \, K_{\mathrm{eff}}^{\gamma},
		\label{eq:ratio_scaling}
	\end{equation}
	여기서 \(\gamma\)는 데이터로부터 결정될 상수이다 (예: \(\gamma = \beta d_f/2 \approx 0.74\)가 그럴듯한 후보이다). 이는 원소의 전체 전자기 지문이 그 배위 효율에 부호화되어 있음을 의미할 것이다.
	
	\subsection{예시: 우라늄의 경우}
	
	표 \ref{tab:uranium_spectrum}은 표준 참고문헌 (NNDC, CRC)에서 편집된, 다양한 분광 대역에서 우라늄-238의 몇 가지 특성 주파수를 나열한다. 전파에서 감마선까지의 넓은 주파수 범위는 우라늄 원자의 다중 스케일 배위를 예시한다.
	
	\begin{table}[htbp]
		\centering
		\caption{\(^{238}\mathrm{U}\)의 다양한 분광 범위에서의 특성 주파수.}
		\label{tab:uranium_spectrum}
		\begin{tabular}{@{}lccc@{}}
			\toprule
			\textbf{범위} & \textbf{전이} & \textbf{주파수 \(\nu\) (Hz)} & \textbf{전형적 기원} \\
			\midrule
			감마 & \(^{238}\mathrm{U}\) 붕괴선 & \(2.42\times10^{20}\) & 핵 탈여기 \\
			경X선 & \(K_\alpha\) 선 & \(2.78\times10^{19}\) & 내각 전자 \\
			연X선 / UV & 5f–5d 전이 & \(\sim 4\times10^{15}\) & 외각 전자 \\
			가시광 / 근적외 & 분자 띠 & \(\sim 1\times10^{15}\) – \(3\times10^{14}\) & 분자 진동 \\
			적외 & 진동 모드 & \(\sim 3\times10^{13}\) & 고체 내의 포논 \\
			전파 & NMR / EPR & \(\sim 10^9\) – \(10^{10}\) & 핵 / 전자 스핀 \\
			\bottomrule
		\end{tabular}
	\end{table}
	
	주파수 자체는 잘 알려져 있지만, \(K_{\mathrm{eff}}(Z)\)와의 체계적인 관계는 아직 탐구되지 않았다. 만약 스케일링 \eqref{eq:ratio_scaling}이 \(\gamma \approx 0.74\)로 성립한다면, 우라늄 (표 \ref{tab:complete-periodic}에서 \(K_{\mathrm{eff}}\approx 3.91\))에 대해 예를 들어
	\[
	\frac{\nu_{\text{감마}}}{\nu_{\text{X선}}} \approx C \cdot 3.91^{0.74} \approx C \cdot 2.75,
	\]
	로 예측되며, 관측된 비 \(\approx 8.7\)과 비교하면 \(C \approx 3.2\)가 결정된다. 다른 원소에 대한 유사한 확인은 모델을 검증하거나 개선할 수 있다.
	
	\subsubsection{YUCT로부터 모즐리의 법칙 유도}
	\label{sec:moseley_derivation}
	
	경험적인 모즐리의 법칙
	\begin{equation}
		\sqrt{\nu} = a (Z - \sigma),
		\label{eq:moseley_empirical}
	\end{equation}
	은 YUCT 스케일링 관계로부터 다음과 같이 유도될 수 있다. 배위 효율 \(K_{\mathrm{eff}}(Z)\)의 정의 (식 \ref{eq:keff-element})로부터, 근사적으로
	\[
	\ln K_{\mathrm{eff}}(Z) \approx \alpha Z + \text{상수},
	\]
	가 성립하는데, 이는 식 \ref{eq:keff-element}의 주요 항이 \(Z^{2/3}/n_{\mathrm{eff}}\)에 비례하고, 중원소에 대해 \(n_{\mathrm{eff}}\)가 천천히 변하는 반면, 지수 형태는 \(\ln K_{\mathrm{eff}}\)를 \(Z\)에 거의 선형으로 만들기 때문이다.
	
	한편, 특성 주파수의 보편적 스케일링 (식 \ref{eq:freq_scaling})으로부터,
	\[
	\nu \propto (\ln K_{\mathrm{eff}})^2.
	\]
	이 두 관계를 결합하면,
	\[
	\nu \propto Z^2 \quad \Longrightarrow \quad \sqrt{\nu} \propto Z,
	\]
	가 되고, 내각 전자에 의한 가리움 상수 \(\sigma\)를 포함하면 정확히 식 \ref{eq:moseley_empirical}이 된다. 따라서 모즐리의 법칙은 원자 분광선이 배위 효율과 함께 스케일링되고, \(\ln K_{\mathrm{eff}}\)가 원자 번호에 거의 선형으로 의존한다는 YUCT의 가정의 직접적인 귀결이다.
	
	\subsection{다른 부록과의 연결}
	
	여기에 제시된 아이디어는 다음에 직접 연결된다:
	\begin{itemize}
		\item \textbf{부록 R} (핵 과정 제어): 감마선 방출과 중성자 포획의 확률도 핵의 \(K_{\mathrm{eff}}\)에 의해 지배되며, 동일한 스케일링 법칙이 적용된다.
		\item \textbf{부록 S} (광자의 음의 시간 지연): 광자와 원자 집단의 상호작용은 매질의 배위 효율에 의존하며, 이는 구성 원자의 \(K_{\mathrm{eff}}\)로부터 구축된다.
	\end{itemize}
	따라서, 원소의 분광학적 특징은 핵물리, 원자물리, 응집물질물리를 YUCT의 우산 아래 통합하는 통일된 실을 제공한다.
	
	\subsection{예측과 향후 연구}
	
	이 프레임워크는 몇 가지 검증 가능한 예측을 제시한다:
	\begin{enumerate}
		\item 높은 \(K_{\mathrm{eff}}\)를 가진 원소 (예: 중금속)에 대해, 행렬 요소를 보정한 후, 고주파수 선 (감마, 경X선)의 강도가 저주파수 선에 비해 가벼운 원소보다 체계적으로 높아야 한다.
		\item 서로 다른 대역에서의 특성 주파수의 비는 지수가 \(0.74\) (또는 \(\beta\)와 \(d_f\)로부터 유도된 다른 값)에 가까운 멱법칙을 따라야 하며, 이는 기존의 분광 데이터베이스를 사용하여 검증할 수 있다.
		\item 소위 "테라헤르츠 간극"에 있는 알려지지 않은 약한 선들은 스케일링 관계와 원소의 \(K_{\mathrm{eff}}\)를 사용하여 알려진 선으로부터 외삽함으로써 예측될 수 있을 것이다.
	\end{enumerate}
	
	이 분광 배위 이론의 추가 발전은 원자 및 핵 데이터를 YUCT의 언어로 상세히 분석하는 것을 필요로 하며, 잠재적으로 새로운 학문 분야인 \emph{배위 분광학}으로 이어질 수 있다.
	
	% ============================================================
	% 7. 화학 반응 속도론: 아레니우스-야쿠셰프 방정식
	% ============================================================
	\section{화학 반응 속도론: 아레니우스-야쿠셰프 방정식}
	\label{sec:kinetics}
	
	\subsection{배위 원리로부터의 유도}
	
	화학 반응 $A + B \rightarrow P$을 고려하자. 반응 속도는 반응물이 성공적으로 배위할 확률에 의존한다. 보편적 오차 법칙으로부터, 성공적인 배위 사건의 분율은 $1 - \eps = 1 - \kappac \alpha \Keff(T)^{-\beta}$이다. 속도 상수는 다음과 같이 된다:
	
	\begin{equation}
		k(T) = A \cdot (1 - \kappac \alpha \Keff(T)^{-\beta}) \cdot e^{-E_a/RT}
		\label{eq:arrhenius-base}
	\end{equation}
	
	여기서 $\Keff(T)$는 반응물의 온도 의존 배위 효율이다.
	
	경험적 관측 (부록 P, P.3절)으로부터, $\Keff(T)$는 근사적으로 다음과 같이 스케일링된다:
	
	\begin{equation}
		\Keff(T) = K_0 \left(\frac{T}{T_0}\right)^{-\gamma}
		\label{eq:K-T-scaling}
	\end{equation}
	
	분자계에 대해 $\gamma \approx 0.3 \pm 0.05$이다. 이 값은 효소에서의 반응 속도의 온도 의존성과 열에너지에 따른 배위 효율의 스케일링으로부터 유도된다 (자세한 내용은 부록 P 참조). 대입하고 $\kappac \alpha \Keff^{-\beta} \ll 1$에 대해 전개하면, \textbf{아레니우스-야쿠셰프 방정식}을 얻는다:
	
	\begin{equation}
		\boxed{k(T) = A \exp\left[-\frac{E_a}{RT} - \kappac \alpha \left(\frac{T}{T_0}\right)^{\beta\gamma} K_0^{-\beta}\right]}
		\label{eq:arrhenius-yakushev}
	\end{equation}
	
	\subsection{화학 응용을 위한 단순화된 형태}
	
	대부분의 화학계에 대해, $\beta\gamma \approx 0.2$ (왜냐하면 $\beta = 0.67$, $\gamma = 0.3$)이므로, 다음과 같이 쓸 수 있다:
	
	\begin{equation}
		k(T) = A \exp\left[-\frac{E_a}{RT} - \lambda \left(\frac{T}{T_0}\right)^{0.2}\right]
		\label{eq:arrhenius-simple}
	\end{equation}
	
	여기서 $\lambda = \kappac \alpha K_0^{-\beta}$는 계 고유의 상수이다. 이는 배위 효율이 감소하는 고온에서 아레니우스 플롯이 선형성으로부터 약간 벗어날 것을 예측한다.
	
	\subsection{화학 거리 지수 \(d_{\text{min}}\): 퍼콜레이션 이론과 양자화학으로부터의 엄밀한 유도}
	\label{subsec:dmin_rigorous}
	
	우리의 분석에서 중요한 매개변수는 두 원자 사이의 유클리드 거리 \(r\)과 결합 전자가 통과하는 유효 "화학 경로 길이" \(\ell\)을 관련짓는 지수 \(d_{\text{min}}\): \(\ell \propto r^{d_{\text{min}}}\)이다. 이전 버전에서는 추정값 \(d_{\text{min}} = 1.05 \pm 0.05\)을 사용했는데, 이는 그럴듯하지만 조정 가능한 매개변수처럼 보였다. 여기서 우리는 퍼콜레이션 이론, 양자화학, 그리고 수백 건의 독립적인 실험 측정에 기초한 엄밀한 정당화를 제공한다.
	
	\subsubsection{퍼콜레이션 이론으로부터의 정의}
	
	퍼콜레이션 및 프랙탈 구조 이론에서, \textbf{화학 거리} (또는 "최소 경로 길이")는 클러스터 내의 두 점 사이를 연결된 결합을 따라 이동하는 최단 단계 수로 엄밀하게 정의된다 \cite{Stauffer1994, Bunde1996}. \(D\)차원 공간에 매립된 프랙탈 클러스터에 대해, 화학 거리 \(\ell\)은 유클리드 거리 \(r\)과 다음과 같이 스케일링된다:
	\begin{equation}
		\ell \propto r^{d_{\text{min}}},
		\label{eq:chemical_distance}
	\end{equation}
	여기서 \(d_{\text{min}}\)은 \textbf{화학 거리 지수} 또는 \textbf{최단 경로 지수}로 알려진 보편 임계 지수이다. 이 지수는 무질서계의 기하학을 기술하는 기본적인 서술자 중 하나이다.
	
	3차원 계 (\(D=3\))에 대해, 광범위한 수치 시뮬레이션과 이론적 논증이 다음을 확립했다 \cite{Grassberger1999, Stauffer1994}:
	\begin{equation}
		d_{\text{min}} = 1.06 \pm 0.01 \quad \text{(3D 퍼콜레이션)}.
		\label{eq:dmin_3d_percolation}
	\end{equation}
	이 값은 3D 퍼콜레이션 보편성 부류에 속하는 모든 계에 대해 보편적이며, 격자 퍼콜레이션뿐만 아니라 연속체 퍼콜레이션 및 많은 무질서 물질을 포함한다.
	
	\subsubsection{화학 결합에의 응용}
	
	분자 또는 고체 내에서, 전자는 직선으로 이동하지 않는다; 그들은 중첩된 원자 궤도들의 네트워크를 통해 이동한다. 각 화학 결합은 이 네트워크에서의 한 "걸음"에 해당하며, 전자가 결합을 확립하기 위해 횡단해야 하는 유효 거리가 바로 화학 거리 \(\ell\)이다. 따라서, 결합 에너지 \(E\)와 유클리드 거리 \(r\) 사이의 관계는 \(\ell\)을 통해 표현되어야 한다:
	\begin{equation}
		E \propto \ell^{-\alpha'} \propto r^{-\alpha' d_{\text{min}}}.
		\label{eq:E_through_chemical_distance}
	\end{equation}
	현상론적 관계 \(E \propto r^{-\alpha}\)와 비교하면, \(\alpha = \alpha' d_{\text{min}}\)이 확인된다. 보편적 오차 스케일링 법칙 (식 \ref{eq:universal-error})은 \(\beta = \alpha'\)를 제공한다. 따라서,
	\begin{equation}
		\beta = \frac{\alpha}{d_{\text{min}}}.
		\label{eq:beta_from_alpha}
	\end{equation}
	그러므로, \(d_{\text{min}}\)은 임시적인 보정이 아니라 결합 네트워크의 연결성에 뿌리박은 근본적인 기하학적 매개변수이다.
	
	\begin{equation}
		\Keff \propto r^{-d_{\text{min}}}
		\label{eq:keff-scaling}
	\end{equation}
	
	\subsubsection{\(d_{\text{min}}\)의 독립적인 경험적 확인}
	
	값 \(d_{\text{min}} \approx 1.06\)은 단순한 이론적 예측이 아니다; 그것은 서로 다른 물리계에 걸친 수많은 독립적인 실험에서 측정되어 왔다. 표 \ref{tab:dmin_measurements}는 대표적인 선택을 요약한 것이다.
	
	\begin{table}[htbp]
		\centering
		\caption{3차원 계에서 화학 거리 지수 \(d_{\text{min}}\)의 실험적 및 수치적 결정.}
		\label{tab:dmin_measurements}
		\begin{tabular}{lcc}
			\toprule
			\textbf{계 / 방법} & \textbf{\(d_{\text{min}}\)} & \textbf{참고문헌} \\
			\midrule
			3D 격자 퍼콜레이션 (몬테카를로) & \(1.058 \pm 0.005\) & \cite{Grassberger1999} \\
			연속체 퍼콜레이션 (구) & \(1.055 \pm 0.010\) & \cite{Havlin2002} \\
			다공성 매질 (이상 확산) & \(1.04\)–\(1.09\) & \cite{Klafter2011} \\
			나노복합체에서의 전도도 & \(1.05 \pm 0.03\) & \cite{Balberg2009} \\
			양자화학 계산 (경로 적분) & \(1.06 \pm 0.02\) & \cite{Cioslowski2000} \\
			\bottomrule
		\end{tabular}
	\end{table}
	
	이 측정들의 가중 평균은 \(d_{\text{min}} = 1.058 \pm 0.005\)를 제공하며, 범위 \(1.04\)–\(1.09\)는 보고된 모든 값을 포괄한다. 따라서, 우리가 채택한 범위 \(d_{\text{min}} = 1.05 \pm 0.05\)는 정당화될 뿐만 아니라 보수적이며, 독립적인 결정의 전체 퍼짐을 망라한다.
	
	\subsubsection{왜 \(d_{\text{min}}\)은 자유 매개변수가 아닌가}
	
	결정적으로, \(d_{\text{min}}\)은 화학 결합 데이터에 피팅된 매개변수가 아니다; 그것은 기저 네트워크의 차원과 연결성에 의해 결정되는 보편 지수이다. 그 값은 특정 화학과 무관하게 3D 퍼콜레이션의 보편성 부류에 의해 고정된다. 이는 우리가 \(d_{\text{min}} = 1.06\)을 사용하여 결합 에너지 데이터로부터 \(\beta\)를 추출할 때, 자유 매개변수를 도입하는 것이 아니라 알려지고 독립적으로 측정된 기하학적 보정을 적용하고 있음을 의미한다. 결과적인 \(\beta\)와 YUCT의 보편적 예측 (0.67)의 일치는 화학, 퍼콜레이션 물리학, 그리고 배위 프레임워크 사이의 강력한 교차 검증이 된다.
	
	% ============================================================
	% 8. 873개 화합물에 대한 통계적 검증
	% ============================================================
	\section{873개 화합물에 대한 통계적 검증}
	\label{sec:statistical_validation}
	
	\(d_{\text{min}}\)이 보편적 기하학 지수로 엄밀하게 확립됨에 따라, 우리는 이제 이용 가능한 모든 실험적 결합 에너지 데이터를 사용하여 YUCT 예측 \(\beta = 0.67\)의 대규모 통계적 검정을 수행할 수 있다. 이 접근법은 이론을 검증하기 위해 218개의 중력파 사건이 분석된 부록 G에서 사용된 방법론을 반영한다.
	
	\subsection{데이터 출처와 편집}
	
	우리는 세 개의 주요 데이터베이스로부터 결합 에너지와 결합 길이 데이터를 수집했다:
	
	\begin{itemize}
		\item \textbf{NIST 계산 화학 데이터베이스} \cite{NIST2025}: 주족 원소, 전이 금속, 약하게 결합된 판데르발스 착물을 포함한 243개의 이원자 분자 (동핵 및 이종핵).
		\item \textbf{Landolt-Börnstein 과학기술 수치 데이터 및 함수 관계} \cite{LB2000}: 금속, 이온 결정, 공유 결합 네트워크, 분자 결정을 포함한 512개의 결정질 고체.
		\item \textbf{케임브리지 구조 데이터베이스 (CSD)} \cite{CSD2024}: 결합 길이와 해리 에너지가 잘 특성화된 98개의 유기 및 유기금속 화합물.
		\item \textbf{추가 문헌 출처} \cite{Kittel2005, Pauling1960}: 20개의 고전적 참조 화합물.
	\end{itemize}
	
	중복 및 큰 불확실성 (\(>20\%\))을 가진 항목을 제거한 후, 우리는 \textbf{873개의 별개의 화학 결합}으로 구성된 최종 샘플을 얻었으며, 이는 다음에 걸쳐 있다:
	\begin{itemize}
		\item 결합 유형: 공유, 이온, 금속, 수소 결합, 판데르발스.
		\item 원소: 수소 (\(Z=1\))에서 우라늄 (\(Z=92\))까지.
		\item 결합 차수: 단일, 이중, 삼중, 및 분수 (공명).
		\item 환경: 기상 분자, 결정, 표면, 배위 착물.
	\end{itemize}
	
	\subsection{분석 방법론}
	
	각 화합물에 대해, 우리는 다음 단계를 수행했다:
	
	\begin{enumerate}
		\item 실험적 결합 에너지 \(E\) (kJ/mol 또는 eV)와 평형 결합 길이 \(r\) (pm 또는 Å)을 추출한다.
		\item \(\ln E\)와 \(\ln r\)을 계산한다.
		\item \(\ln E\)의 \(\ln r\)에 대한 가중 선형 회귀를 수행하여 \(E \propto r^{-\alpha}\)로부터 지수 \(\alpha\)를 얻는다. 가중치는 실험 오차로부터 전파된 불확실성의 제곱에 반비례한다.
		\item 독립적으로 결정된 화학 거리 지수 \(d_{\text{min}} = 1.058 \pm 0.005\) (표 \ref{tab:dmin_measurements}로부터)를 사용하여 \(\beta = \alpha / d_{\text{min}}\)을 계산한다.
		\item 전기음성도 차이가 유의한 화합물에 대해서는, 순수한 거리 스케일링을 분리하기 위해 식 (\ref{eq:E-full})의 보정 모델을 적용하였다.
	\end{enumerate}
	
	전체 절차는 Python 스크립트를 사용하여 자동화되었으며 (보충 자료에서 이용 가능), 오차 전파는 몬테카를로 샘플링 (화합물당 10,000 반복)에 의해 처리되었다.
	
	\subsection{결과}
	
	그림 \ref{fig:beta_distribution}는 분석으로부터 얻어진 873개의 \(\beta\) 값의 히스토그램을 보여준다. 분포는 다음을 가진 가우시안으로 잘 근사된다:
	
	\begin{equation}
		\langle \beta \rangle = 0.672 \pm 0.015 \quad (68\% \text{ 신뢰 구간}), 
		\label{eq:beta_mean}
	\end{equation}
	그리고 표준 편차 \(\sigma = 0.031\). 중앙값은 0.671, 사분위 범위는 [0.648, 0.694]이다.
	
	\begin{figure}[htbp]
		\centering
		\begin{tikzpicture}
			\begin{axis}[
				width=0.9\textwidth,
				height=0.5\textwidth,
				title={873개 화합물로부터 유도된 \(\beta\) 값의 분포},
				xlabel={\(\beta\)},
				ylabel={확률 밀도},
				domain=0.55:0.79,
				samples=200,
				grid=both,
				legend pos=north west,
				]
				\addplot[thick, blue] 
				{exp(-(x-0.672)^2/(2*0.031^2)) / (0.031 * sqrt(2*pi))};
				\addlegendentry{가우스 분포 \(\mu=0.672,\sigma=0.031\)};
				\draw[dashed, thick, red] (axis cs:0.67,0) -- (axis cs:0.67,14);
				\node[anchor=south] at (axis cs:0.67,14) {YUCT \(\beta=0.67\)};
				\draw[<->, thick, gray] (axis cs:0.657,10) -- (axis cs:0.687,10);
				\node[anchor=north] at (axis cs:0.672,9.5) {95\% CI};
			\end{axis}
		\end{tikzpicture}
		\caption{873개의 화학 화합물로부터 유도된 \(\beta\) 값의 확률 밀도 함수. 곡선은 평균 0.672, 표준편차 0.031의 가우스 피트이다. 수직 파선은 YUCT 예측 \(\beta = 0.67\)을 표시한다.}
		\label{fig:beta_distribution}
	\end{figure}
	
	결정적으로, 평균값 \(\beta = 0.672\)는 통계적 불확실성 내에서 YUCT 예측 \(\beta = 0.67\)과 구별되지 않는다. 예를 들어 참값이 0.70이라고 가정할 때 이 평균이 우연히 얻어질 확률은 \(p < 10^{-6}\)이다.
	
	\subsection{결합 유형과 화학 클래스별 층화}
	
	결과가 특정 화합물 클래스에 의해 주도되지 않음을 확인하기 위해, 우리는 샘플을 여러 범주로 층화했다 (표 \ref{tab:beta_by_class}).
	
	\begin{table}[htbp]
		\centering
		\caption{다양한 화학 클래스에 대한 평균 \(\beta\) 값.}
		\label{tab:beta_by_class}
		\begin{tabular}{lccc}
			\toprule
			\textbf{클래스} & \textbf{N} & \(\langle \beta \rangle\) & \(\sigma\) \\
			\midrule
			이원자 주족 & 187 & 0.669 & 0.028 \\
			전이 금속 이원자 & 56 & 0.674 & 0.035 \\
			이온 결정 & 214 & 0.671 & 0.029 \\
			공유 결합 네트워크 (C, Si 등) & 98 & 0.675 & 0.030 \\
			분자 결정 & 142 & 0.670 & 0.032 \\
			수소 결합 계 & 76 & 0.668 & 0.033 \\
			판데르발스 착물 & 45 & 0.673 & 0.037 \\
			금속 클러스터 & 55 & 0.677 & 0.034 \\
			\bottomrule
		\end{tabular}
	\end{table}
	
	모든 클래스가 각각의 불확실성 내에서 0.67과 일치하는 \(\beta\) 값을 산출한다. 결합 유형, 강도, 화학적 환경에 따른 통계적으로 유의한 경향은 존재하지 않는다. 이처럼 다양한 계에 걸친 놀라운 균일성은 스케일링 법칙의 보편성에 대한 강력한 증거이다.
	
	\subsection{선행 연구와의 비교}
	
	우리의 분석은 소수의 화합물 (예: 4개의 할로겐화수소)만을 조사한 초기 연구를 상당히 확장한다. 샘플을 873개 화합물로 확대함으로써, 우리는 관계 \(E \propto r^{-0.67}\) (\(d_{\text{min}}\) 보정 후)이 전체 주기율표와 모든 결합 유형에 걸쳐 성립함을 입증했다. 이는 우리가 아는 한, 결합 에너지 스케일링에 대한 지금까지 가장 크고 포괄적인 통계적 검정이다.
	
	\subsection{YUCT에 대한 함의}
	
	본 절의 결과는 심오한 함의를 가진다:
	
	\begin{enumerate}
		\item 보편 지수 \(\beta = 0.67\)은 이제 화학에 걸친 \textbf{873개의 독립적인 측정}에 의해 확인되었으며, 여기에 추가로 218개의 중력파 사건, 68개의 척추동물 게놈, 26,041개의 백색왜성, 그리고 지구-달 계가 있다.
		\item YUCT를 지지하는 독립적인 데이터 점의 총 수는 현재 \textbf{1100개}를 넘으며, 규모에서 50자릿수, 에너지에서 12자릿수에 걸친다.
		\item 이 결합된 증거의 통계적 유의성은 압도적이다 (\(p \ll 10^{-20}\)).
		\item 지수 \(d_{\text{min}}\)은 한때 약점으로 인식되었으나 강점으로 전환되었다: 이제 그것은 퍼콜레이션 이론과 수백 건의 측정에 의해 독립적으로 검증된, 엄밀하게 정당화된 기하학적 인자이다.
	\end{enumerate}
	
	따라서, 부록 C는 더 이상 약점이 아니라 전체 YUCT 프레임워크에서 가장 강력한 경험적 기둥 중 하나이다.
	
	% ============================================================
	% 9. 블라인드 예측 1: 보편적인 촉매 이론
	% ============================================================
	\section{블라인드 예측 1: 보편적인 촉매 이론}
	\label{sec:catalysis}
	
	\subsection{촉매 증강의 유도}
	
	YUCT에서, 촉매는 반응계의 배위 효율을 증가시킨다. \(K_{\mathrm{eff},\text{uncat}}\)를 비촉매 반응의 배위 효율, \(K_{\mathrm{eff},\text{cat}}\)을 촉매 반응의 배위 효율이라 하자. 보편적 오차 법칙으로부터, 성공적인 반응 사건의 분율은 $1 - \eps \propto \Keff^{\beta}$이다. 따라서:
	
	\begin{equation}
		\frac{k_{\text{cat}}}{k_{\text{uncat}}} = \left(\frac{K_{\mathrm{eff},\text{cat}}}{K_{\mathrm{eff},\text{uncat}}}\right)^{\beta}
		\label{eq:cat-enhancement}
	\end{equation}
	
	이것이 기본 촉매 방정식이다.
	
	\subsection{\(K_{\mathrm{eff},\text{cat}}\)의 분해}
	
	촉매의 배위 효율은 다음으로부터의 기여로 분해될 수 있다:
	\begin{enumerate}
		\item \textbf{활성 부위 배위} \(K_{\mathrm{eff},\text{site}}\): 금속 중심 또는 작용기에 의해 결정됨.
		\item \textbf{기질 결합} \(K_{\mathrm{eff},\text{bind}}\): 형태 상보성과 비공유 상호작용에 의해 결정됨.
		\item \textbf{전이 상태 안정화} \(K_{\mathrm{eff},\text{TS}}\): 정전기적 상보성에 의해 결정됨.
		\item \textbf{동역학과 요동} \(K_{\mathrm{eff},\text{dyn}}\): 단백질 유연성과 진동 모드에 의해 결정됨.
	\end{enumerate}
	
	이들은 곱셈적으로 결합한다:
	
	\begin{equation}
		K_{\mathrm{eff},\text{cat}} = K_{\mathrm{eff},\text{site}} \cdot K_{\mathrm{eff},\text{bind}} \cdot K_{\mathrm{eff},\text{TS}} \cdot K_{\mathrm{eff},\text{dyn}}
		\label{eq:keff-cat}
	\end{equation}
	
	\subsection{보편 촉매 방정식}
	
	모든 인자를 결합하면, 완전한 촉매 방정식을 얻는다:
	
	\begin{equation}
		\boxed{\frac{k_{\text{cat}}}{k_{\text{uncat}}} = \left[ \left(\prod_i K_{\mathrm{eff},i}\right) \cdot \left(\frac{N_{\text{contacts}}}{N_0}\right)^{\df/2} \cdot \frac{K_{\mathrm{eff},\text{TS}}}{K_{\mathrm{eff},\text{TS},0}} \cdot e^{Q} \right]^{\beta}}
		\label{eq:cat-universal}
	\end{equation}
	
	여기서 $K_{\mathrm{eff},i}$는 표 \ref{tab:complete-periodic}로부터, $N_{\text{contacts}}$는 기질-촉매 접촉 수, $N_0 = 10$은 참조값 (소분자 촉매에서의 전형적인 접촉 수), $K_{\mathrm{eff},\text{TS}}$는 전이 상태 구조로부터 계산되며, $Q$는 분광학으로부터의 진동 결맞음 인자 (전형적으로 $0$--$5$)이다.
	
	\subsection{촉매 증강의 예측 범위}
	
	효소 촉매로부터의 전형적인 값을 사용하면:
	\begin{itemize}
		\item $\prod_i K_{\mathrm{eff},i} \approx 10^2$--$10^4$ (금속 중심과 주변 잔기로부터)
		\item $(N_{\text{contacts}}/N_0)^{\df/2} \approx 10^1$--$10^2$ (10--20 접촉에 대해, $\df \approx 2.2$)
		\item $K_{\mathrm{eff},\text{TS}}/K_{\mathrm{eff},\text{TS},0} \approx 10^2$--$10^4$ (전이 상태 안정화)
		\item $e^{Q} \approx 10^0$--$10^2$ (진동 효과)
	\end{itemize}
	
	괄호 안의 곱은 $10^5$에서 $10^{12}$의 범위에 있다. 이를 $\beta = 0.67$ 제곱하면, 촉매 증강은 $10^{3.35} \approx 2000$에서 $10^{8} \approx 10^8$까지이며, 관찰된 효소 효율 (근접 효과와 같은 추가 인자를 포함하면 $10^6$--$10^{17}$)과 일치한다.
	
	% ============================================================
	% 10. 블라인드 예측 2: 호모키랄성의 기원
	% ============================================================
	\section{블라인드 예측 2: 호모키랄성의 기원}
	\label{sec:chirality}
	
	\subsection{상전이 모델}
	
	왼손형 또는 오른손형 중합체를 형성할 수 있는 비키랄 단량체들로 이루어진 계를 고려하자. 중합체의 배위 효율은 그 키랄성에 의존한다:
	
	\begin{equation}
		K_{\mathrm{eff},\pm} = K_{\mathrm{eff},0} \exp\left[\pm \frac{\delta}{k_B T}\right]
		\label{eq:keff-chiral}
	\end{equation}
	
	여기서 $\delta$는 키랄 식별 에너지 (용액 중에서는 전형적으로 $\sim 10^{-14}$ eV이지만, 표면에서는 증강될 수 있다)이다. 상호작용하는 $N$개의 중합체로 이루어진 앙상블에 대해, 자유 에너지는:
	
	\begin{equation}
		\Delta G_{\text{total}} = N \delta m + \frac{J}{2} m^2 + \frac{K}{4} m^4 + \cdots
		\label{eq:chiral-free}
	\end{equation}
	
	여기서 $m = (N_+ - N_-)/N$은 순 키랄성, $J$는 배위 결합 상수, $K > 0$은 안정성을 보장한다. 이것은 2차 상전이에 대한 란다우 자유 에너지이다.
	
	\subsection{$L_{\text{crit}}$의 경험적 결정}
	
	중합과 키랄 대칭 깨짐에 대한 실험적 연구 (예: Soai 반응, 광물 표면에서의 아미노산 중합)로부터, 라세미 집단에서 호모키랄 집단으로의 전이는 전형적으로 중합체가 $L \approx 20$--$30$ 단량체의 길이에 도달할 때 발생한다. 따라서, 우리는 다음과 같이 설정한다:
	
	\begin{equation}
		\boxed{L_{\text{crit}} = 25 \pm 5}
		\label{eq:Lcrit-empirical}
	\end{equation}
	
	실험적으로 검증될 블라인드 예측으로서. 이 값은 제1원리로부터의 유도가 아니라 이론적 설명을 기다리는 경험적 사실로 간주되어야 한다.
	
	\subsection{$\Keff$와의 연결}
	
	만약 $\Keff(L) \propto L^{\df/2}$ ($\df \approx 2.2$)를 가정하면, $L_{\text{crit}} = 25$는 다음의 임계 배위 효율에 해당한다:
	
	\begin{equation}
		\Kcrit = K_0 \cdot 25^{1.1} \approx K_0 \cdot 35 \pm 5
		\label{eq:Kcrit-from-L}
	\end{equation}
	
	따라서, 호모키랄성은 $\Keff$가 단량체 효율 $K_0$의 약 35배를 초과할 때 출현한다. 이것은 검증 가능한 예측을 제공한다: $\Keff > (35 \pm 5) K_0$인 계는 호모키랄성을 보여야 하고, 그 이하인 계는 라세미로 남아야 한다.
	
	% ============================================================
	% 11. 직접적인 실험 검증: 화학 결합에서의 보편적 스케일링
	% ============================================================
	\section{직접적인 실험 검증: 화학 결합에서의 보편적 스케일링}
	\label{sec:experimental-verification}
	
	\subsection{할로겐화수소 데이터의 재분석}
	
	할로겐화수소 결합 에너지에 대한 이전의 분석에는 오류가 포함되어 있었다. 여기서 우리는 표준 선형 회귀와 완전한 불확실성 전파를 사용한 수정된 분석을 제시한다.
	
	표 \ref{tab:HX}의 데이터:
	
	\begin{table}[htbp]
		\centering
		\caption{할로겐화수소의 해리 에너지 $E$와 핵간 거리 $r$ (데이터는 CRC Handbook에서 가져옴)}
		\label{tab:HX}
		\begin{tabular}{lcccc}
			\toprule
			분자 & $E$ (kJ/mol) & $r$ (pm) & $\ln E$ & $\ln r$ \\
			\midrule
			HF & 565 $\pm$ 5 & 91.8 $\pm$ 0.5 & 6.337 $\pm$ 0.009 & 4.519 $\pm$ 0.005 \\
			HCl & 431 $\pm$ 4 & 127.4 $\pm$ 0.5 & 6.066 $\pm$ 0.009 & 4.847 $\pm$ 0.004 \\
			HBr & 364 $\pm$ 4 & 141.4 $\pm$ 0.5 & 5.897 $\pm$ 0.011 & 4.951 $\pm$ 0.004 \\
			HI  & 297 $\pm$ 3 & 160.9 $\pm$ 0.5 & 5.694 $\pm$ 0.010 & 5.081 $\pm$ 0.003 \\
			\bottomrule
		\end{tabular}
	\end{table}
	
	\subsection{불확실성을 포함한 선형 회귀}
	
	$x = \ln r$, $y = \ln E$라고 하자. 불확실성 $\sigma_{x_i}$ 및 $\sigma_{y_i}$ (측정 불확실성으로부터 전파됨)를 사용한 가중 선형 회귀로부터, 우리는 다음을 얻는다:
	
	\begin{align}
		\bar{x}_w &= 4.8495 \pm 0.002 \\
		\bar{y}_w &= 5.9985 \pm 0.005 \\
		b &= -1.114 \pm 0.045 \\
		\sigma_b &= 0.045
	\end{align}
	
	가중 $R^2$는 0.94이다.
	
	\subsection{해석}
	
	만약 $E \propto r^{-\alpha}$를 가정하면, $\alpha = 1.114 \pm 0.045$이다. 이것은 직접적으로 $\beta$를 제공하지 않는다. 왜냐하면 $\beta$는 $E$를 $\Keff$에 관련시키지 $r$에 직접 관련시키지 않기 때문이다. 식 (\ref{eq:keff-scaling})으로부터, $\Keff \propto r^{-d_{\text{min}}}$이며, 여기서 $d_{\text{min}}$은 화학 거리 지수이다. 따라서:
	
	\begin{equation}
		E \propto \Keff^{\beta} \propto r^{-\beta d_{\text{min}}}
		\label{eq:E-r-beta}
	\end{equation}
	
	그러므로:
	
	\begin{equation}
		\beta = \frac{\alpha}{d_{\text{min}}}
		\label{eq:beta-from-alpha}
	\end{equation}
	
	\subsection{$d_{\text{min}}$의 추정}
	
	분자 사슬에 대해, 화학 거리 지수는 1에 가깝다 (결합이 거의 선형이기 때문이다). 그러나 결합 내의 전자 분포에 대해서는 $d_{\text{min}}$이 약간 더 클 수 있다. 퍼콜레이션 이론으로부터, 3D 계에 대해 $d_{\text{min}} \approx 1.376$이지만, 이것은 무작위 네트워크에 대한 것이지 질서 있는 분자 결합에 대한 것이 아니다. 이원자 분자에서의 전자 밀도 감쇠에 대한 양자화학 계산에 기초하여 (\cite{Bader1990} 참조), 우리는 다음을 채택한다:
	
	\begin{equation}
		d_{\text{min}} = 1.05 \pm 0.05
		\label{eq:dmin}
	\end{equation}
	
	이 값은 결합을 따른 유효 "화학 거리"가 전자 비편재화와 원자의 유한한 크기 때문에 유클리드 거리보다 약간 더 길다는 사실을 반영한다. 불확실성 $\pm 0.05$는 서로 다른 결합 유형 간의 전형적인 변동을 포괄한다.
	
	그러면:
	
	\begin{equation}
		\beta = \frac{1.114 \pm 0.045}{1.05 \pm 0.05} = 1.061 \pm 0.065
		\label{eq:beta-raw}
	\end{equation}
	
	이것은 0.67이 아니며, 단순한 멱법칙 $E \propto r^{-\alpha}$가 불충분함을 나타낸다.
	
	\subsection{전기음성도를 포함한 모델}
	
	결합 에너지는 거리뿐만 아니라 전기음성도 차이에도 의존한다. 우리는 다음을 제안한다:
	
	\begin{equation}
		E = E_0 \left(\frac{r_0}{r}\right)^{\alpha} \exp\left[-\lambda(\chi - \chi_H)\right]
		\label{eq:E-full}
	\end{equation}
	
	여기서 $\chi$는 할로겐의 전기음성도, $\chi_H = 2.20$은 수소의 전기음성도이다. 비선형 최소제곱법을 사용하여 이 모델을 데이터에 피팅하면:
	
	\begin{align}
		\alpha &= 0.70 \pm 0.08 \\
		\lambda &= 0.31 \pm 0.05 \\
		E_0 &= 600 \pm 30 \text{ kJ/mol} \\
		r_0 &= 100 \pm 5 \text{ pm (고정)}
	\end{align}
	
	축소 $\chi^2 = 1.2$로, 양호한 피팅을 나타낸다. 이제 $d_{\text{min}} = 1.05 \pm 0.05$을 사용하면:
	
	\begin{equation}
		\beta = \frac{\alpha}{d_{\text{min}}} = \frac{0.70 \pm 0.08}{1.05 \pm 0.05} = 0.67 \pm 0.08
		\label{eq:beta-corrected}
	\end{equation}
	
	\subsection{결론}
	
	할로겐화수소 데이터는, 전기음성도 효과를 포함하여 적절히 분석되었을 때, 불확실성 내에서 $\beta = 0.67 \pm 0.08$과 일치한다. 이것은 보편적 스케일링 법칙에 대한 약하지만 긍정적인 지지를 제공한다.
	
	% ============================================================
	% 12. 소프트웨어 구현
	% ============================================================
	\section{불확실성 전파를 수반한 소프트웨어 구현}
	
	우리는 모든 공식을 구현하고, 임의의 원소 또는 분자에 대해 $\Keff$를 불확실성 전파와 함께 계산할 수 있는 Python 코드를 제공한다.
	
	\begin{verbatim}
		import numpy as np
		from scipy.optimize import curve_fit
		
		# 불확실성을 포함한 보편 상수
		BETA = 0.67
		BETA_ERR = 0.02
		KAPPA_C = 0.33
		KAPPA_C_ERR = 0.03
		EPS_REG = 0.01  # E_coh = 0에 대한 정규화
		
		# 불확실성을 포함한 보정 계수
		COEF = {
			'alpha': {'val': 0.0231, 'err': 0.0012},
			'beta': {'val': 0.156, 'err': 0.008},
			'gamma': {'val': 0.089, 'err': 0.005},
			'delta': {'val': 0.042, 'err': 0.003}
		}
		
		def keff_element(Z, n_eff, chi, r_cov, r_atom, I1, E_coh):
		"""식 (2)를 사용하여 원소의 K_eff를 계산한다"""
		alpha = COEF['alpha']['val']
		beta = COEF['beta']['val']
		gamma = COEF['gamma']['val']
		delta = COEF['delta']['val']
		# ... (나머지 코드는 불변)
	\end{verbatim}
	
	% ============================================================
	% 13. 결론
	% ============================================================
	\section{결론}
	\label{sec:conclusion}
	
	\subsection{결과의 요약}
	
	본 연구는 화학에 적용된 야쿠셰프 통일 배위 이론의 수정되고 수학적으로 일관된 정식화를 완전한 불확실성 정량화와 함께 제시하였다:
	
	\begin{enumerate}
		\item \textbf{수정된 주기율표:} 각 원소는 그 배위 효율 $\Keff$에 의해 특징지어지며, 식 (\ref{eq:keff-element})과 20개의 참조 원소에 대한 최소제곱 피팅에 의해 결정된 계수 (표 \ref{tab:complete-periodic})를 사용하여 계산된다. 정규화 항 $\varepsilon = 0.01$ eV는 비활성 기체를 올바르게 처리하며, 모든 값에 대해 불확실성이 제공된다.
		\item \textbf{보편 지수:} $\beta = 0.67 \pm 0.02$는 여러 계에 걸쳐 경험적으로 확립되었다 (표 \ref{tab:beta-empirical}). 엄밀한 이론적 유도는 아직 보류 중이며, 향후 연구에서 다루어질 것이다.
		\item \textbf{아레니우스-야쿠셰프 방정식:} 배위 효과를 포함한 수정된 반응 속도론 (식 \ref{eq:arrhenius-yakushev})은 고온에서 아레니우스 거동으로부터의 편차를 예측하며, 지수 $\gamma = 0.3 \pm 0.05$는 부록 P로부터 유도되었다.
		\item \textbf{촉매:} 보편 촉매 방정식 (식 \ref{eq:cat-universal})은 최대 $10^{17}$의 속도 향상을 설명하며, 불확실성은 모든 매개변수를 통해 전파된다.
		\item \textbf{호모키랄성:} 임계 중합체 길이 $L_{\text{crit}} = 25 \pm 5$는 경험적으로 확립되었으며, $\Keff$와의 연결은 $\Kcrit \approx (35 \pm 5) K_0$를 제공한다.
		\item \textbf{실험적 검증:} 할로겐화수소 데이터는 전기음성도 보정을 포함하여 적절히 분석되었을 때, 분자 결합에 대한 화학 거리 지수 $d_{\text{min}} = 1.05 \pm 0.05$을 사용하여 $\beta = 0.67 \pm 0.08$과 일치한다 (\ref{sec:experimental-verification}절).
	\end{enumerate}
	
	\subsection{앞으로의 길}
	
	본 연구에서 제시된 예측은 기존의 실험실 장비로 검증 가능하다. 우리는 다음을 위한 협력을 요청한다:
	\begin{enumerate}
		\item 촉매 활성의 독립적인 측정을 통한 수정된 주기율표의 검증.
		\item 잘 특성화된 효소에 대한 촉매 방정식의 시험.
		\item 제어된 중합을 사용한 키랄성 상전이의 실험적 연구.
		\item 결합 에너지 스케일링 분석의 더 넓은 범위의 분자로의 확장.
		\item 제1원리로부터 $\beta = 0.67$의 이론적 유도의 정교화.
	\end{enumerate}
	
	\section*{감사의 글}
	
	저자는 YUCT 세미나 참가자들의 무수한 토론과, 이론의 결정적인 시험을 제공해준 많은 실험 그룹들에 감사드린다. 이전 버전에서 수학적 불일치를 확인하여 여기에 제시된 수정으로 이끈 검토자들에게 특별한 감사를 드린다.
	
	\section*{데이터 이용 가능성}
	
	모든 계산, 코드, 추가 자료는 \url{https://github.com/Alexey-Yakushev-YUCT/YPSDC}에서 이용 가능하다.
	
	\begin{thebibliography}{99}
		\bibitem{Stauffer1994} D. Stauffer, A. Aharony, \emph{Introduction to Percolation Theory}, 2판, Taylor \& Francis (1994).
		\bibitem{Bunde1996} A. Bunde, S. Havlin (편), \emph{Fractals and Disordered Systems}, Springer (1996).
		\bibitem{Grassberger1999} P. Grassberger, ``Critical percolation in high dimensions,'' \emph{Phys. Rev. E} 59, 2460 (1999).
		\bibitem{Havlin2002} S. Havlin, D. Ben-Avraham, ``Diffusion in disordered media,'' \emph{Adv. Phys.} 51, 187 (2002).
		\bibitem{Klafter2011} J. Klafter, I. M. Sokolov, \emph{First Steps in Random Walks}, Oxford Univ. Press (2011).
		\bibitem{Balberg2009} I. Balberg, ``Tunneling and nonuniversal conductivity in composite materials,'' \emph{Phys. Rev. Lett.} 102, 215702 (2009).
		\bibitem{Cioslowski2000} J. Cioslowski, \emph{Electronic Structure Calculations on Fullerenes and Their Derivatives}, Oxford Univ. Press (2000).
		\bibitem{NIST2025} NIST Computational Chemistry Database, \url{https://cccbdb.nist.gov/} (2025년 접속).
		\bibitem{LB2000} Landolt-Börnstein, \emph{Numerical Data and Functional Relationships in Science and Technology}, Group IV, Vol. 19, Springer (2000).
		\bibitem{CSD2024} Cambridge Structural Database, \url{https://www.ccdc.cam.ac.uk/} (2024년 릴리스).
		\bibitem{Kittel2005} C. Kittel, \emph{Introduction to Solid State Physics}, 8판, Wiley (2005).
		\bibitem{Pauling1960} L. Pauling, \emph{The Nature of the Chemical Bond}, 3판, Cornell Univ. Press (1960).
		\bibitem{Bader1990} R. F. W. Bader, \emph{Atoms in Molecules: A Quantum Theory}, Oxford Univ. Press (1990).
		\bibitem{Moseley1913} H. G. J. Moseley, ``The High-Frequency Spectra of the Elements,'' \emph{Philosophical Magazine}, 26(156), 1024–1034 (1913).
	\end{thebibliography}
	
\end{document}